% ============================================================================
% 05_reinforcement_learning/02_actor_critics.tex
% Section body only: include immediately after the corresponding \section
% command in the master book file. No document-level commands are included.
% ============================================================================

The preceding discussion established reinforcement learning as an optimization
problem in which an agent improves a policy through interaction with an
environment and feedback from a reward signal. The next question is how the
learning algorithm can determine which individual decisions deserve
reinforcement. An observed return evaluates the consequences of a sequence of
decisions, whereas policy optimization must update the probability of the
individual actions that produced that sequence. This distinction becomes
especially important in autoregressive language generation: a response is
formed through many token-level decisions, but its quality may only be
observable after a substantial part, or all, of the response has been
generated.

Policy-gradient methods provide a direct way to optimize a parameterized
stochastic policy. Their practical effectiveness, however, depends on the
quality of the signal used to evaluate sampled actions. Returns can vary
substantially across trajectories that begin in similar states, and this
variation can make policy-gradient estimates noisy. The actor--critic
architecture addresses this problem by separating policy representation from
value estimation. The actor represents the policy that selects actions; the
critic estimates the expected future return associated with the states in
which those actions are selected. The critic's estimate supplies a baseline
for constructing advantages, allowing the actor to learn from the relative
quality of an action rather than from its raw return alone.

This separation is central to understanding Proximal Policy Optimization
(PPO), which combines a policy objective with a learned value function
\parencite{schulman2017ppo}. It is also important for understanding why later
language-model methods investigate alternatives to a separately trained
critic. The present section develops the actor--critic formulation from its
probabilistic and statistical foundations, derives the principal estimators
used by PPO, and then maps these components onto autoregressive language
models.

\subsection{Sequential decision-making and the policy represented by the actor}
\label{subsec:ac-policy-actor}

A sequential decision problem can be represented by a Markov decision process
(MDP),
\begin{equation}
    \mathcal{M}
    =
    (\mathcal{S},\mathcal{A},P,r,\gamma),
    \label{eq:ac-mdp}
\end{equation}
where \(\mathcal{S}\) is the state space, \(\mathcal{A}\) is the action space,
\(P(s'\mid s,a)\) is the transition distribution, \(r\) is the reward
function, and \(\gamma\in[0,1]\) is the discount factor. The Markov property
means that, conditional on the current state and action, the distribution of
the next state does not depend on the full history preceding that state. The
state is therefore intended to retain the information needed to characterize
the future evolution of the decision process.

At time \(t\), the agent occupies state \(s_t\), samples an action \(a_t\),
receives reward \(r_t\), and moves to state \(s_{t+1}\). A trajectory over a
finite horizon \(T\) can be written as
\begin{equation}
    \tau
    =
    (s_0,a_0,r_0,s_1,a_1,r_1,\ldots,s_{T-1},a_{T-1},r_{T-1},s_T).
    \label{eq:ac-trajectory}
\end{equation}
The discounted return from time \(t\) is
\begin{equation}
    G_t
    =
    \sum_{k=t}^{T-1}
    \gamma^{k-t}r_k.
    \label{eq:ac-return}
\end{equation}
The discount factor determines how strongly rewards further in the future
contribute to the return. When \(\gamma=0\), only the immediate reward
contributes. As \(\gamma\) approaches one, later rewards receive greater
weight. For an episodic finite-horizon problem, the sum terminates at the
episode boundary.

The agent's behavior is described by a policy. In the stochastic,
parameterized setting considered by policy-gradient methods, the policy is
\begin{equation}
    \pi_\theta(a_t\mid s_t),
    \label{eq:ac-policy}
\end{equation}
where \(\theta\) denotes the trainable policy parameters. For a fixed state,
\(\pi_\theta(\cdot\mid s)\) is a probability distribution over the available
actions. The policy objective is the expected return under trajectories
generated by that policy:
\begin{equation}
    J(\theta)
    =
    \mathbb{E}_{\tau\sim\pi_\theta}[G_0],
    \qquad
    \theta^\star
    =
    \arg\max_\theta J(\theta).
    \label{eq:ac-objective}
\end{equation}

The actor is the parameterized function that implements this policy. In a
neural implementation, it maps a state representation to the parameters of
an action distribution. For a discrete action space with \(K\) actions, let
\(\mathbf{z}_\theta(s)\in\mathbb{R}^{K}\) denote the actor's logits. The
probability assigned to action \(k\) is obtained through the softmax map:
\begin{equation}
    \pi_\theta(a=k\mid s)
    =
    \frac{\exp(z_{\theta,k}(s))}
    {\displaystyle\sum_{j=1}^{K}\exp(z_{\theta,j}(s))}.
    \label{eq:ac-softmax}
\end{equation}
The resulting vector belongs to the probability simplex: every component is
nonnegative and the components sum to one. The dimensionality of the actor's
output is therefore determined by the size of the discrete action space.

For a differentiable stochastic policy, the policy-gradient estimator used
in PPO has the form
\begin{equation}
    \widehat{\mathbf{g}}
    =
    \widehat{\mathbb{E}}_t
    \left[
        \nabla_\theta
        \log\pi_\theta(a_t\mid s_t)
        \widehat{A}_t
    \right],
    \label{eq:ac-policy-gradient}
\end{equation}
where \(\widehat{\mathbb{E}}_t\) denotes an empirical average over sampled
time steps and \(\widehat{A}_t\) is an estimate of the advantage. The
log-probability gradient describes how the probability of the sampled action
changes as the policy parameters change. The advantage supplies the
evaluation signal: a positive estimate encourages an increase in the
probability of the sampled action under gradient ascent, whereas a negative
estimate encourages a decrease. The policy-gradient paper formulates the
corresponding surrogate objective as the empirical expectation of
\(\log\pi_\theta(a_t\mid s_t)\widehat{A}_t\)
\parencite{schulman2017ppo}.

For a language model, the state at generation step \(t\) consists of the
prompt and the tokens generated so far. If \(x\) denotes the prompt and
\(y_{<t}=(y_1,\ldots,y_{t-1})\) denotes the generated prefix, then
\begin{equation}
    s_t=(x,y_{<t}),
    \qquad
    a_t=y_t,
    \qquad
    \pi_\theta(a_t\mid s_t)
    =
    \pi_\theta(y_t\mid x,y_{<t}).
    \label{eq:ac-llm-state-action}
\end{equation}
The action is the next token selected from the vocabulary. If the vocabulary
contains \(|\mathcal{V}|\) tokens, the actor's language-model output
distribution has \(|\mathcal{V}|\) components. The probability of a complete
response \(y=(y_1,\ldots,y_T)\) follows from the autoregressive
factorization:
\begin{equation}
    \pi_\theta(y\mid x)
    =
    \prod_{t=1}^{T}
    \pi_\theta(y_t\mid x,y_{<t}),
    \qquad
    \log\pi_\theta(y\mid x)
    =
    \sum_{t=1}^{T}
    \log\pi_\theta(y_t\mid x,y_{<t}).
    \label{eq:ac-autoregressive-factorization}
\end{equation}
Thus, language-model policy optimization is a sequence of token-level
stochastic decisions, even when the reward is assigned to the completed
response.

\subsection{Value functions, the critic, and the advantage}
\label{subsec:ac-value-critic}

The actor determines which actions are sampled, but the policy-gradient
estimator also requires an estimate of how well a sampled action performed.
The value-function formulation provides this evaluation. The state-value
function of a policy \(\pi\) is the expected discounted return conditional
on the current state:
\begin{equation}
    V^\pi(s_t)
    =
    \mathbb{E}_{\pi}
    \left[
        G_t\mid s_t
    \right].
    \label{eq:ac-state-value}
\end{equation}
It is an expectation over possible future trajectories, not the return of
one particular sampled trajectory. Different trajectories can begin from
the same state and produce different returns because the policy is
stochastic and the transition process may be stochastic.

The action-value function conditions on both the current state and the
action selected at that state:
\begin{equation}
    Q^\pi(s_t,a_t)
    =
    \mathbb{E}_{\pi}
    \left[
        G_t\mid s_t,a_t
    \right].
    \label{eq:ac-action-value}
\end{equation}
The state-value function is the policy-weighted average of the action-value
function:
\begin{equation}
    V^\pi(s)
    =
    \sum_{a\in\mathcal{A}}
    \pi(a\mid s)Q^\pi(s,a),
    \label{eq:ac-vq-relation}
\end{equation}
with an integral replacing the summation for a continuous action space.
The corresponding Bellman expectation equation is
\begin{equation}
    V^\pi(s)
    =
    \sum_{a\in\mathcal{A}}\pi(a\mid s)
    \sum_{s'\in\mathcal{S}}P(s'\mid s,a)
    \left[
        r(s,a,s')+\gamma V^\pi(s')
    \right].
    \label{eq:ac-bellman}
\end{equation}
This equation expresses the expected value of a state in terms of the
immediate reward and the discounted value of the next state, averaged over
the policy's actions and the environment's transitions.

The advantage function compares the value of a particular action with the
average value of the current policy at the same state:
\begin{equation}
    A^\pi(s,a)
    =
    Q^\pi(s,a)-V^\pi(s).
    \label{eq:ac-advantage}
\end{equation}
A positive advantage means that the action has a greater expected return
than the policy's average action at that state. A negative advantage means
that it has a lower expected return. The comparison is state-conditional:
an absolute return that is large in one state may be poor in another state
if the latter state normally leads to even larger returns.

The advantage has zero conditional expectation under the policy:
\begin{align}
    \mathbb{E}_{a\sim\pi(\cdot\mid s)}
    [A^\pi(s,a)]
    &=
    \mathbb{E}_{a\sim\pi(\cdot\mid s)}
    [Q^\pi(s,a)-V^\pi(s)]
    \nonumber\\
    &=
    \mathbb{E}_{a\sim\pi(\cdot\mid s)}[Q^\pi(s,a)]
    -V^\pi(s)
    \nonumber\\
    &=
    V^\pi(s)-V^\pi(s)
    =
    0.
    \label{eq:ac-zero-mean-advantage}
\end{align}
This property is central to the statistical role of the critic.

In a practical actor--critic system, the exact value function is generally
not available. The critic therefore learns a parameterized approximation:
\begin{equation}
    V_\phi(s)\approx V^\pi(s),
    \label{eq:ac-critic-approximation}
\end{equation}
where \(\phi\) denotes the critic parameters. Unlike the actor, whose output
is a probability distribution over actions, the critic produces a scalar
estimate of expected future return. If the state representation has
dimension \(d_s\), a conventional critic can be viewed as a function
\(V_\phi:\mathbb{R}^{d_s}\rightarrow\mathbb{R}\).

The critic is trained by regression against a target return or a return
estimate. For a target \(\widehat{V}^{\,\mathrm{target}}_t\), the squared
value-function error is
\begin{equation}
    \mathcal{L}_{V}(\phi)
    =
    \frac{1}{2}
    \mathbb{E}
    \left[
        \left(
            V_\phi(s_t)-\widehat{V}^{\,\mathrm{target}}_t
        \right)^2
    \right].
    \label{eq:ac-value-loss}
\end{equation}
When the empirical return \(G_t\) is used as the target, this becomes
\begin{equation}
    \mathcal{L}_{V}(\phi)
    =
    \frac{1}{2}
    \mathbb{E}
    \left[
        \left(V_\phi(s_t)-G_t\right)^2
    \right].
    \label{eq:ac-monte-carlo-value-loss}
\end{equation}
The factor \(1/2\) is a convenient scaling factor. Differentiating the
objective gives
\begin{equation}
    \nabla_\phi\mathcal{L}_{V}(\phi)
    =
    \mathbb{E}
    \left[
        \left(V_\phi(s_t)-G_t\right)
        \nabla_\phi V_\phi(s_t)
    \right].
    \label{eq:ac-value-gradient}
\end{equation}
Under squared-error minimization, the optimal prediction for a given state
is the conditional expectation of the target. The critic is consequently
a regression model for expected future return, rather than a predictor of
the exact future trajectory.

The value function also provides a baseline for policy-gradient estimation.
To see why subtracting a state-dependent baseline does not change the
expected policy gradient, consider an arbitrary function \(b(s_t)\) that
does not depend on the sampled action once the state is fixed. Its
conditional contribution to the gradient is
\begin{align}
    &\mathbb{E}_{a_t\sim\pi_\theta(\cdot\mid s_t)}
    \left[
        \nabla_\theta\log\pi_\theta(a_t\mid s_t)b(s_t)
    \right]
    \nonumber\\
    &\quad=
    b(s_t)
    \sum_{a_t}
    \pi_\theta(a_t\mid s_t)
    \nabla_\theta\log\pi_\theta(a_t\mid s_t)
    \nonumber\\
    &\quad=
    b(s_t)
    \nabla_\theta
    \sum_{a_t}\pi_\theta(a_t\mid s_t)
    =
    b(s_t)\nabla_\theta 1
    =
    0.
    \label{eq:ac-baseline-identity}
\end{align}
The identity uses
\(\pi_\theta(a\mid s)\nabla_\theta\log\pi_\theta(a\mid s)
=\nabla_\theta\pi_\theta(a\mid s)\) and the normalization of the policy
distribution. Therefore, replacing \(G_t\) by \(G_t-b(s_t)\) preserves the
expected policy gradient. Choosing \(b(s_t)=V_\phi(s_t)\) gives the
practical return-minus-value advantage estimate
\begin{equation}
    \widehat{A}_t
    =
    G_t-V_\phi(s_t).
    \label{eq:ac-monte-carlo-advantage}
\end{equation}
The baseline does not need to be exact for the baseline identity to hold;
it must, however, be a function of the state and not of the sampled action
conditional on that state. The accuracy of the critic still matters for
the variance of the estimator and, when bootstrapped targets are used, for
the bias of the resulting advantage estimate.

\subsection{Temporal-difference estimation and generalized advantage estimation}
\label{subsec:ac-td-gae}

The empirical return in \cref{eq:ac-monte-carlo-value-loss} uses the
rewards observed over the subsequent trajectory. This Monte Carlo target
does not require a value prediction for the next state, but it can have
substantial variance because it depends on all subsequent rewards. An
alternative is temporal-difference (TD) estimation, which uses the
immediate reward and the critic's estimate of the next state's value.

The one-step TD residual is
\begin{equation}
    \delta_t
    =
    r_t+\gamma V_\phi(s_{t+1})-V_\phi(s_t).
    \label{eq:ac-td-residual}
\end{equation}
The first term is the observed immediate reward. The second is the
discounted estimate of the next state's value. The final term is the
critic's current estimate. A positive residual indicates that the
immediate reward plus the estimated continuation value exceeds the
critic's current prediction; a negative residual indicates the opposite.
Because the next-state value is supplied by the critic, this estimator
uses bootstrapping.

Generalized Advantage Estimation (GAE), used in the PPO formulation,
combines TD residuals across future time steps
\parencite{schulman2017ppo}. For a finite trajectory segment ending at
time \(T\), its estimator is
\begin{equation}
    \widehat{A}^{\mathrm{GAE}(\gamma,\lambda)}_t
    =
    \sum_{\ell=0}^{T-t-1}
    (\gamma\lambda)^\ell\delta_{t+\ell},
    \qquad
    \delta_{t+\ell}
    =
    r_{t+\ell}
    +\gamma V_\phi(s_{t+\ell+1})
    -V_\phi(s_{t+\ell}),
    \label{eq:ac-gae}
\end{equation}
where \(\lambda\in[0,1]\) controls the weighting of future TD residuals.
The discount factor \(\gamma\) controls the contribution of future rewards,
whereas \(\lambda\) controls the degree to which the estimator relies on
longer sequences of residuals rather than short-horizon bootstrapping.

When \(\lambda=0\), the estimator reduces to the one-step TD residual:
\begin{equation}
    \widehat{A}^{\mathrm{GAE}(\gamma,0)}_t
    =
    \delta_t.
    \label{eq:ac-gae-lambda-zero}
\end{equation}
When \(\lambda=1\), the residuals telescope:
\begin{align}
    \widehat{A}^{\mathrm{GAE}(\gamma,1)}_t
    &=
    \sum_{\ell=0}^{T-t-1}
    \gamma^\ell
    \left[
        r_{t+\ell}
        +\gamma V_\phi(s_{t+\ell+1})
        -V_\phi(s_{t+\ell})
    \right]
    \nonumber\\
    &=
    \sum_{\ell=0}^{T-t-1}\gamma^\ell r_{t+\ell}
    +\gamma^{T-t}V_\phi(s_T)-V_\phi(s_t).
    \label{eq:ac-gae-telescoping}
\end{align}
For a completed episode with terminal value \(V_\phi(s_T)=0\), this is
the Monte Carlo return minus the current value estimate. For a truncated
trajectory segment that has not reached a terminal state, the final
bootstrap term \(\gamma^{T-t}V_\phi(s_T)\) remains. This is why a
finite-horizon implementation can use a value estimate at the end of a
collected segment rather than requiring the complete episode to be
observed.

The parameter \(\lambda\) controls a bias--variance trade-off. Smaller
values place more weight on short-horizon TD residuals and therefore rely
more heavily on the critic's bootstrap predictions. This can reduce
sampling variance, but errors in the critic can then influence the
advantage estimate more directly. Larger values incorporate more sampled
rewards and rely less on intermediate bootstrap predictions, but the
estimate can inherit more of the variance of the sampled trajectory. The
choice of \(\lambda\) therefore balances the statistical properties of
the sampled returns against the approximation quality of the critic.

For example, suppose \(r_t=0\), \(\gamma=1\),
\(V_\phi(s_t)=0.60\), and \(V_\phi(s_{t+1})=0.80\). The first TD residual
is
\begin{equation}
    \delta_t
    =
    0+0.80-0.60
    =
    0.20.
    \label{eq:ac-gae-example-delta}
\end{equation}
If the next residual is \(\delta_{t+1}=-0.10\) and
\(\lambda=0.5\), the two-residual estimate is
\begin{equation}
    \widehat{A}_t
    =
    \delta_t+\gamma\lambda\delta_{t+1}
    =
    0.20+(1)(0.5)(-0.10)
    =
    0.15.
    \label{eq:ac-gae-example-result}
\end{equation}
The positive result indicates that the sampled action performed better
than the critic's current expectation, after incorporating the specified
future residual. The calculation also shows how a later negative residual
can reduce, but need not reverse, the estimated advantage.

\subsection{The actor--critic objective in Proximal Policy Optimization}
\label{subsec:ac-ppo}

The critic supplies an advantage estimate, but it does not by itself
constrain how far the actor moves during optimization. If the actor is
updated repeatedly using the same sampled trajectories, its action
probabilities can move substantially away from those of the policy that
generated the data. PPO addresses this issue with a clipped surrogate
objective \parencite{schulman2017ppo}.

Let \(\theta_{\mathrm{old}}\) denote the policy parameters used to collect
the current batch of trajectories. The probability ratio for a sampled
state--action pair is
\begin{equation}
    \rho_t(\theta)
    =
    \frac{\pi_\theta(a_t\mid s_t)}
    {\pi_{\theta_{\mathrm{old}}}(a_t\mid s_t)}.
    \label{eq:ac-importance-ratio}
\end{equation}
At the beginning of the update, when
\(\theta=\theta_{\mathrm{old}}\), the ratio equals one. Values above one
mean that the current policy assigns a larger probability to the sampled
action than the data-collection policy; values below one mean that it
assigns a smaller probability.

The unclipped surrogate term is
\(\rho_t(\theta)\widehat{A}_t\). PPO modifies it by clipping the ratio to
an interval around one:
\begin{equation}
    \mathcal{L}^{\mathrm{CLIP}}(\theta)
    =
    \widehat{\mathbb{E}}_t
    \left[
        \min
        \left(
            \rho_t(\theta)\widehat{A}_t,\,
            \operatorname{clip}
            \left(
                \rho_t(\theta),1-\epsilon,1+\epsilon
            \right)\widehat{A}_t
        \right)
    \right],
    \label{eq:ac-ppo-clipped-objective}
\end{equation}
where \(\epsilon>0\) is a hyperparameter controlling the clipping interval.
The minimum makes the objective conservative: it removes the incentive
for a probability-ratio change beyond the clipping threshold when that
change would otherwise improve the surrogate objective. It does not impose
a hard guarantee that every policy distribution remains inside a fixed
distance of the old policy; it modifies the optimization objective so that
certain large changes no longer provide additional improvement.

The effect depends on the sign of the advantage. If
\(\widehat{A}_t>0\), increasing the probability of the sampled action
improves the unclipped term, but the clipped term stops providing additional
incentive once the ratio exceeds \(1+\epsilon\). If
\(\widehat{A}_t<0\), decreasing the probability of the sampled action
improves the objective, but the lower clipping threshold prevents further
surrogate improvement once the ratio falls below \(1-\epsilon\).

For a numerical illustration, let \(\epsilon=0.2\),
\(\rho_t(\theta)=1.35\), and \(\widehat{A}_t=0.4\). The unclipped
contribution is \(1.35(0.4)=0.54\), whereas the clipped contribution is
\(1.2(0.4)=0.48\). The minimum is therefore \(0.48\). The objective does
not reward the additional probability increase beyond the upper clipping
threshold.

The critic is trained alongside the actor. Let
\(V_t^{\mathrm{target}}\) denote the return or bootstrapped target used for
value regression. The value-function error is
\begin{equation}
    \mathcal{L}^{\mathrm{VF}}_t(\phi)
    =
    \left(
        V_\phi(s_t)-V_t^{\mathrm{target}}
    \right)^2.
    \label{eq:ac-ppo-value-loss}
\end{equation}
PPO can also include an entropy bonus to discourage premature concentration
of the policy distribution. For a discrete action distribution, its
entropy is
\begin{equation}
    \mathcal{H}(\pi_\theta(\cdot\mid s_t))
    =
    -\sum_{a\in\mathcal{A}}
    \pi_\theta(a\mid s_t)
    \log\pi_\theta(a\mid s_t).
    \label{eq:ac-policy-entropy}
\end{equation}
The combined actor--critic objective, written as an objective to maximize,
can be expressed as
\begin{align}
    \mathcal{J}^{\mathrm{PPO}}(\theta,\phi)
    =
    \widehat{\mathbb{E}}_t
    \Big[
        &\ell^{\mathrm{CLIP}}_t(\theta)
        -c_1\mathcal{L}^{\mathrm{VF}}_t(\phi)
        \nonumber\\
        &+c_2\mathcal{H}(\pi_\theta(\cdot\mid s_t))
    \Big],
    \label{eq:ac-ppo-joint-objective}
\end{align}
where \(\ell^{\mathrm{CLIP}}_t\) is the per-sample term inside the
expectation in \cref{eq:ac-ppo-clipped-objective}, and \(c_1,c_2\) weight
the value error and entropy bonus. The negative sign on the value error
reflects the fact that the combined objective is maximized, whereas the
critic's regression error is minimized.

The actor and critic may use separate parameter sets or share part of a
neural representation. When parameters are shared, the policy surrogate
and value error both contribute gradients to the shared representation.
The PPO paper explicitly describes this combined objective and the
possibility of shared parameters between the policy and value function
\parencite{schulman2017ppo}. The relative weighting of the two losses then
matters because the actor and critic impose different learning signals on
the shared representation.

The resulting training cycle alternates between data collection and
optimization. The current policy generates trajectories; the algorithm
computes reward and value information; advantage estimates are constructed;
and the actor and critic are optimized over minibatches for several epochs.
The policy used for the next collection phase is then updated to the newly
optimized policy. The following algorithm summarizes this actor--critic
cycle in the notation used for language-model training.

\begin{algorithm}[H]
    \caption{PPO-style actor--critic optimization for autoregressive language models}
    \label{alg:ac-ppo-llm}
    \KwIn{Initial policy \(\pi_\theta\), reference policy
    \(\pi_{\mathrm{ref}}\), reward model \(R_\psi\), prompt distribution
    \(\mathcal{D}\), discount \(\gamma\), GAE parameter \(\lambda\),
    clipping parameter \(\epsilon\), and optimization budget}
    \KwOut{Updated policy parameters \(\theta\)}
    Initialize or load the critic \(V_\phi\)\;
    \For{each training iteration until the optimization budget is exhausted}{
        Sample a batch of prompts from \(\mathcal{D}\)\;
        Set \(\theta_{\mathrm{old}}\leftarrow\theta\)\;
        Generate responses from \(\pi_{\theta_{\mathrm{old}}}\)\;
        Evaluate completed responses with \(R_\psi\)\;
        Construct token-level rewards, including the reference-policy penalty\;
        Evaluate the critic on the prompt and each response prefix\;
        Compute TD residuals and GAE advantages\;
        Construct return targets for value-function regression\;
        \For{each optimization epoch and minibatch}{
            Compute the probability ratios using
            \(\pi_\theta/\pi_{\theta_{\mathrm{old}}}\)\;
            Update the actor using the clipped surrogate objective\;
            Update the critic using the value-function regression objective\;
        }
    }
\end{algorithm}

The actor in this algorithm is the language model's policy. This should not
be confused with the parallel data-collection actors sometimes used in
conventional PPO implementations: those are rollout workers that execute
the policy to collect trajectories. Schulman et al. describe the
alternation between parallel trajectory collection and multiple epochs of
minibatch optimization in their actor--critic-style PPO algorithm
\parencite{schulman2017ppo}.

\subsection{Mapping the actor--critic architecture onto language models}
\label{subsec:ac-llm-mapping}

The MDP representation of language generation makes the roles of the actor
and critic explicit. A prompt \(x\) initiates an episode, a partial response
\(x\circ y_{<t}\) is the state, and the next token \(y_t\) is the action.
After the token is appended, the next state is
\(x\circ y_{\leq t}\). Ivison et al. describe language generation in this
way when detailing PPO for preference-based language-model training
\parencite{ivison2024unpacking}.

At the actor's output interface, a language-model hidden representation
\(\mathbf{h}_t\in\mathbb{R}^{d_{\mathrm{model}}}\) is mapped to a vector of
vocabulary logits
\(\mathbf{z}_t\in\mathbb{R}^{|\mathcal{V}|}\). A softmax converts these
logits into the token distribution. The critic, in contrast, maps the
representation of the same prompt--prefix state to a scalar:
\begin{equation}
    V_\phi(x\circ y_{<t})\in\mathbb{R}.
    \label{eq:ac-llm-value}
\end{equation}
In an implementation with a scalar regression head, the critic can use a
representation of dimension \(d_{\mathrm{model}}\) and map it to one
real-valued prediction. The regression head can be evaluated at each
response-token position, so the critic estimates the expected return of
incomplete responses rather than assigning a value only to the completed
sequence. Ivison et al. report that their value model uses an architecture
similar to the reward model, with a regression head that can be applied at
response-token positions \parencite{ivison2024unpacking}.

The reward model and the reference model have different roles from the
critic. A reward model \(R_\psi(x,y)\) assigns a scalar score to a completed
response. In preference-based training, it is learned from comparisons
between responses. The reference policy \(\pi_{\mathrm{ref}}\) is a fixed
policy against which the optimized actor is regularized. The critic
\(V_\phi(x\circ y_{<t})\) estimates expected future return from an
intermediate state. Although the reward model and critic both produce
scalar outputs, they do not estimate the same quantity: one scores a
response, while the other predicts the expected return from a partial
response under the active policy.

In the PPO formulation described by Ivison et al., the sequence-level
reward is applied at the final response token, while a token-level
regularization term penalizes divergence from the reference policy. A
corresponding token-level reward can be written as
\begin{equation}
    r_t
    =
    -\beta
    \log
    \frac{
        \pi_\theta(y_t\mid x,y_{<t})
    }{
        \pi_{\mathrm{ref}}(y_t\mid x,y_{<t})
    }
    +
    \mathbb{I}[t=T]R_\psi(x,y),
    \label{eq:ac-llm-token-reward}
\end{equation}
where \(\beta\) controls the strength of the reference-policy penalty and
\(\mathbb{I}[t=T]\) is one at the final response token and zero otherwise.
The discounted return from a response prefix is then
\begin{equation}
    G_t
    =
    \sum_{k=t}^{T}
    \gamma^{k-t}r_k,
    \label{eq:ac-llm-return}
\end{equation}
with the terminal indexing understood according to the response boundary.
The critic estimates the expected value of this return from the current
prompt--prefix state. Its estimate is used to form token-level advantages,
which then guide the actor's update.

\cref{fig:ac-architecture} summarizes the information flow. The actor
generates a token from the current state. The resulting response is
evaluated through the reward process. In parallel, the critic estimates
the value of the current state. The reward and value estimate jointly
determine the advantage signal, which is used to update the actor. The
diagram is a conceptual synthesis of the PPO actor--critic formulation
and its language-model interpretation
\parencite{schulman2017ppo,ivison2024unpacking}.

\begin{figure}[htbp]
    \centering
    \resizebox{0.95\textwidth}{!}{% Figure source for 02_actor_critics.tex.
% Uses the standard LaTeX picture environment and requires no additional
% drawing package. Include this file with \input from the project root.
\setlength{\unitlength}{1mm}
\begin{picture}(155,80)
    % Main policy-generation path.
    \put(0,51){\framebox(28,14){\shortstack{State\\\(s_t=(x,y_{<t})\)}}}
    \put(42,51){\framebox(30,14){\shortstack{Actor\\\(\pi_\theta\)}}}
    \put(87,51){\framebox(20,14){\shortstack{Token\\\(y_t\)}}}
    \put(120,51){\framebox(28,14){\shortstack{Reward\\\(r_t\)}}}
    \put(28,58){\vector(1,0){14}}
    \put(72,58){\vector(1,0){15}}
    \put(107,58){\vector(1,0){13}}

    % Value-estimation and advantage path.
    \put(42,20){\framebox(30,14){\shortstack{Critic\\\(V_\phi(s_t)\)}}}
    \put(87,20){\framebox(28,14){\shortstack{Advantage\\\(\widehat A_t\)}}}
    \put(120,20){\framebox(32,14){\shortstack{Actor update\\\(\theta\leftarrow\theta'\)}}}
    \put(28,51){\vector(1,-1){18}}
    \put(72,27){\vector(1,0){15}}
    \put(120,51){\vector(-1,-1){17}}
    \put(115,27){\vector(1,0){5}}

    % Feedback from the policy update to the actor.
    \put(152,27){\line(1,0){1}}
    \put(153,27){\line(0,1){48}}
    \put(57,75){\line(1,0){96}}
    \put(57,75){\vector(0,-1){10}}
\end{picture}
}
    \caption{Information flow in an actor--critic language-model training
    step. The actor maps the prompt and generated prefix to a token
    distribution, while the critic estimates the value of that prefix.
    Reward information and the value estimate are combined to construct
    an advantage used in the policy update. Conceptual synthesis based on
    \textcite{schulman2017ppo} and \textcite{ivison2024unpacking}.}
    \label{fig:ac-architecture}
\end{figure}

The value model is trained by minimizing a regression error against
empirical or bootstrapped returns. Ivison et al. describe the value loss
as a mean-squared error between the predicted value of each incomplete
response and its empirical return. They also report implementation
choices in which the value model is initialized from the reward model,
while other implementations initialize it from a supervised-fine-tuned
or base model and replace the language-model output head with a regression
head \parencite{ivison2024unpacking}. These are implementation choices,
not defining properties of the actor--critic architecture.

\subsection{Computational limitations and the transition beyond a learned critic}
\label{subsec:ac-limitations-transition}

The critic provides a principled baseline for variance-reduced
policy-gradient estimation, but it introduces a second learned model into
the training system. In large language models, a value model with a
parameter scale comparable to that of the policy can impose substantial
memory and computational costs. The cost is not limited to storing the
additional parameters: the critic must also be evaluated on response
prefixes and trained through a value-regression objective. These
requirements become more consequential when responses are long and the
critic must produce estimates at many token positions.

A further difficulty arises when the reward model supplies a sequence-level
score primarily at the end of a response. The critic must then learn to
predict the expected eventual outcome from intermediate prefixes, even
though the decisive reward signal may not be available until generation
finishes. Shao et al. identify both the computational burden of a
value model comparable in size to the policy and the difficulty of
training accurate token-level values when the reward is concentrated at
the final token \parencite{shao2024deepseekmath}.

These limitations motivate methods that estimate a baseline without
training a separate value model. One example is Group Relative Policy
Optimization (GRPO), introduced by Shao et al. as a PPO variant. For a
prompt \(q\), the method samples several responses
\(\{y_1,\ldots,y_G\}\) from the old policy and evaluates their rewards
\(\{R_1,\ldots,R_G\}\). A group mean can be formed as
\begin{equation}
    \overline{R}
    =
    \frac{1}{G}\sum_{i=1}^{G}R_i,
    \qquad
    \sigma_R
    =
    \sqrt{
        \frac{1}{G}
        \sum_{i=1}^{G}(R_i-\overline{R})^2
    }.
    \label{eq:ac-group-statistics}
\end{equation}
The normalized relative reward used as a group-based advantage is
\begin{equation}
    \widehat{A}_i
    =
    \frac{R_i-\overline{R}}{\sigma_R},
    \label{eq:ac-group-advantage}
\end{equation}
when the group standard deviation is nonzero. The response-level estimate
can be associated with the tokens of the corresponding response. This
construction replaces the learned state-value baseline with a baseline
calculated from the rewards of multiple responses to the same prompt.
Shao et al. use this group-relative comparison to remove the additional
value-function approximation in GRPO \parencite{shao2024deepseekmath}.

A group-relative baseline and a learned critic are not identical
estimators. The critic predicts a state-dependent expected return for an
individual prompt--prefix state. A group baseline compares completed
responses generated for the same prompt. The latter avoids training an
additional value model, but it does not provide the same token-level
state-value approximation. If all responses in a group receive the same
reward, then \(\sigma_R=0\) and the normalized expression in
\cref{eq:ac-group-advantage} is undefined; such a group contains no
within-group reward variation from which to construct a relative ranking.
The group-based construction therefore changes the statistical source of
the baseline rather than eliminating the need to estimate relative
performance.

The actor--critic architecture establishes how a learned value function
can convert delayed trajectory outcomes into state-conditioned advantage
estimates, and how those estimates can be combined with a constrained
policy update. Its limitations also clarify why later language-model
algorithms investigate alternative estimators and reward structures. The
next section turns from the architecture used to estimate value toward
the reward signal itself, with particular attention to reinforcement
learning in which outcomes can be checked against verifiable task
criteria.
